\documentclass[11pt,a4paper]{article}

\usepackage[utf8]{inputenc}
\usepackage{amsmath, amssymb, amsthm}
\usepackage{graphicx}
\usepackage{enumitem}
\usepackage[margin=1in]{geometry}

% One command per assignment detail, so the header is edited in a
% single place rather than scattered through the document.
\newcommand{\course}{Course 101}
\newcommand{\assignment}{{name}}
\newcommand{\student}{Your Name}

\newtheorem*{solution}{Solution}

\title{\assignment}
\author{\student \\ \small \course}
\date{\today}

\begin{document}

\maketitle

\section*{Problem 1}

State the problem here, then answer it.

\begin{solution}
    Work through the argument in full. Displayed steps are easier to
    follow — and easier to mark — than a single dense paragraph:
    \begin{align}
        (a + b)^2 &= (a + b)(a + b) \\
                  &= a^2 + 2ab + b^2.
    \end{align}
\end{solution}

\section*{Problem 2}

A problem with parts:

\begin{enumerate}[label=(\alph*)]
    \item The first part.
    \item The second part.
\end{enumerate}

\begin{solution}
    \begin{enumerate}[label=(\alph*)]
        \item The answer to the first part.
        \item The answer to the second part.
    \end{enumerate}
\end{solution}

\section*{Problem 3}

\begin{solution}
    Remember to say what a result means, not just that it holds.
\end{solution}

\end{document}
